\section{Cold War Context and the Space Race Begins}

\subsection{Superpower rivalry after 1945}
After the Second World War, the United States and the Soviet Union emerged as rival superpowers whose competition shaped global politics for decades. Their contest was ideological, military, and economic, but it was also deeply symbolic: each side sought visible proof that its system could out-build and out-think the other. Space, still an untouched frontier, offered exactly that kind of stage.

The rivalry hardened quickly. Berlin, Korea, and the nuclear arms race left little room for trust, and reliable intelligence on the other side was scarce and prized. In that climate, technological milestones doubled as strategic signals: a successful launch told allies and adversaries alike that a nation could deliver a warhead across continents, so the line between civilian achievement and military deterrence was deliberately blurred.

Walter McDougall's classic political history argues that the space race forced both governments to become technocratic states: to plan, fund, and direct science on a scale that peacetime democracies and command economies alike had never attempted. Spaceflight was therefore never a private adventure; it was state policy, organized around budgets, five-year plans, and election cycles.

The competition also reshaped daily life far from any launch pad. American classrooms gained new science curricula, Soviet engineering institutes expanded, and both publics learned to read rocket launches as scoreboard entries in a global contest. The vocabulary of the era, from missile gap to space age, captured how thoroughly the rivalry had fused technology with identity.

Orbit offered something practical as well as symbolic: a vantage point no treaty could close. Reconnaissance satellites promised photographs of territory that aircraft could reach only at great risk, which is one reason both governments tolerated each other's early satellites instead of treating an overflight as an act of war. The quiet legal precedent set in those first years, that space belongs to no one, proved as consequential as any single launch.

Neither government ever spent purely for show. Every public milestone rested on a foundation of classified work in guidance, reentry, and tracking whose budgets dwarfed the civilian headlines, and historians still debate where the propaganda program ended and the weapons program began. That ambiguity was the point: a nation that could orbit a satellite on demand had proven things no parade could.

Prestige, then, was never a side effect of the space race; it was a central goal. Governments invested enormous resources because being first carried weight in capitals, newspapers, and living rooms around the world. The Moon eventually became the ultimate prize precisely because it was difficult, expensive, and visible to everyone who looked up at night. \cite{siddiqi2010} \cite{mcdougall1985} \cite{cadbury2006}

\subsection{Rockets, missiles, and prestige}
The rockets that carried satellites and cosmonauts were direct descendants of wartime weapons. Germany's V-2, the first large liquid-fueled ballistic missile, proved that controlled rocket flight at the edge of space was possible, and at the war's end both superpowers raced to recruit its engineers and seize its hardware. The expertise gathered in 1945 became the seed of two parallel national programs.

Through the 1950s, missile development and space ambition advanced together. The same boosters designed to deliver nuclear warheads could, with modest changes, loft scientific payloads, so investment in one capability quietly strengthened the other. This dual use meant that every rocket test was watched closely abroad, and that failures were as politically costly as they were technically instructive.

Behind the hardware stood two very different engineering cultures. The Soviet effort concentrated authority in secret design bureaus led by a chief designer whose word settled disputes, while the American program distributed work across military services, private contractors, and laboratories that competed for funding. Each arrangement produced brilliance and waste in its own way, and each left fingerprints on the vehicles it built.

Propellant chemistry, guidance electronics, and lightweight structures all matured in this decade under military budgets. By the time political leaders demanded spectacular space achievements, the underlying technologies were close enough to ready that engineers could promise results within years rather than decades, a confidence that would have been impossible in 1950.

The first launch vehicles made the lineage explicit. The rocket that carried Sputnik was a lightly modified R-7 intercontinental missile, and America's early satellites rode boosters drawn from the Redstone and Jupiter missile programs. Civilian space agencies would later commission dedicated launchers, but the race began on hardware designed to carry warheads, a fact every government involved understood perfectly well.

Prestige tied the whole enterprise together. A reliable, powerful rocket signaled industrial maturity, scientific depth, and military credibility in a single stroke. Because these signals were read by allies, adversaries, and domestic audiences at once, governments were willing to fund ambitious programs whose payoff was measured as much in influence as in engineering. \cite{harford1997} \cite{logsdon2010}

\begin{table}[htbp]
\centering
\caption{Selected Moon Race milestones}
\label{tab:milestones}
\begin{tabular}{@{}lll@{}}
\toprule
Year & Program & Milestone \\
\midrule
1957 & Sputnik & First artificial satellite \\
1961 & Vostok & First human orbital flight \\
1965 & Gemini & Rendezvous and EVA practice \\
1969 & Apollo 11 & First crewed lunar landing \\
1975 & Apollo--Soyuz & Post-race cooperation \\
\bottomrule
\end{tabular}
\end{table}

\begin{figure}[htbp]
\centering
\includegraphics[width=0.88\linewidth]{figures/ch01_cold_war.jpg}
\caption{Apollo 11 Saturn V on Launch Pad 39A before the Moon landing mission.}
\label{fig:ch_apollo_11_saturn_v_on_launch_pad}
\end{figure}

\bigskip
\begin{otherlanguage}{hebrew}
\subsection*{סיכום הפרק}

פרק זה מציג את רקע המלחמה הקרה ואת ראשית מרוץ החלל. שתי מעצמות-העל הפכו טילים בליסטיים להישג יוקרתי שנועד לתעמולה ולהוכחת עליונות טכנולוגית.

\end{otherlanguage}

\clearpage
